%!TEX root = ../../document/document.tex

% For every chapter create a file in the /content/chapters folder with the name
% chapter-XX.tex where XX is the chapter number  (e.g. 01, 02, 03, ..., 99).


\chapter{Einleitung}\label{c:1}

\section{Problemstellung}\label{c:1_s:problem}
Die Cloud-Technologie hat in den letzten Jahren erheblich an Bedeutung gewonnen. Unternehmen und Organisationen nutzen zunehmend Cloud-Dienste, um ihre \ac{IT}-Infrastruktur zu optimieren. Dabei gibt es verschiedene Modelle, die unterschiedliche Anforderungen und Bedürfnisse abdecken. Die drei Hauptmodelle sind \ac{IaaS}, \ac{PaaS} und \ac{SaaS}. Diese Modelle bieten unterschiedliche Ebenen der Abstraktion und Kontrolle über die zugrunde liegende Infrastruktur.\\
Da diese verschiedenen Systeme immer komplexer werden, ist es umso wichtiger, dass Unternehmen ihre IT-Infrastruktur wiederholbar, effizient und kostengünstig aufbauen können. Hierfür eignet sich das Verwenden von \ac{IaC}-Tools, um die Bereitstellung der Infrastruktur automatisiert wird.

\section{Aufgabenstellung}\label{c:1_s:task}
Die Aufgabe dieser Arbeit besteht darin, die drei Aspekte \ac{IaaS}, \ac{PaaS} und \ac{SaaS} anhand eines Praxisbeispiels verwenden zu können.\\
Hierfür soll für den \ac{SaaS}-Anteil eine Webanwendung entwickelt werden, die es Nutzern ermöglicht, Texte in verschiedene Sprachen zu übersetzen. Diese Webanwendung soll dann auf einer Microsoft Azure VM bereitgestellt werden, um den Aspekt von \ac{IaaS} zu erfüllen. Um letztendlich ebenfalls noch den \ac{PaaS}-Aspekt mit integrieren zu können, soll mit einer Datenbank kommuniziert werden, um bereits getätigte Anfragen zwischenspeichern zu können. Das Deployment von den Services soll mithilfe von Terraform und Ansible automatisiert werden. 
